\chapter{Analysis Pipeline}
\label{chap:analysis-pipeline}

\section{Preprocessing Pipeline}
\label{sec:pipeline-preprocessing}

% TODO: The fixed correction order and why it's fixed: frame sanity check
% -> saturation check -> baseline subtraction -> flat-field division ->
% bad-pixel masking -> optional ROI masking. Frame this as encoded in
% run_preprocessing(), not just documented convention.

\section{Noise Model and Uncertainty Propagation}
\label{sec:pipeline-noise-model}

% TODO: Thompson-Larson-Webb-style noise model combining background_sigma
% (from baseline calibration) and conversion gain; how uncertainty
% propagates through to the fitted dispersion coefficients.

\section{Spatial Dispersion Fitting}
\label{sec:pipeline-dispersion-fitting}

% TODO: x0(lambda) = c0 + c1*lambda, zeta = dx0/dlambda = c1; total
% least-squares fitting (why TLS over OLS -- both axes carry uncertainty).
% See presentation slide 15 for the fit diagram/equations.

\subsection{Optical Design Trade-off: Focal-Length Selection}
\label{sec:pipeline-focal-length-tradeoff}

% TODO: Now that the noise model (Section~\ref{sec:pipeline-noise-model})
% and the TLS fit above exist, revisit the focal lengths stated descriptively
% in Section~\ref{sec:optical-layout} (collimating/focussing doublet pair,
% periscope relay lenses) and show quantitatively how they were chosen to
% minimise the propagated uncertainty in the fitted zeta -- i.e. this is
% the payoff section for the forward reference planted in Chapter 2, not a
% restatement of the optical layout itself.

\section{Multi-Shot Combination}
\label{sec:pipeline-multi-shot}

% TODO: How repeated shots are combined into one result (reduced
% chi-squared as the diagnostic), and the design choice of averaging before
% vs after centroiding -- flag this as something later re-examined as a
% possible systematic in the Discussion chapter.
